\documentclass[11pt,a4paper]{article}

\usepackage{amsmath,amssymb,amsthm,mathtools}
\usepackage[utf8]{inputenc}
\usepackage[margin=1in]{geometry}
\usepackage{hyperref}
\usepackage{enumitem}
\usepackage{booktabs}

% Theorem environments
\theoremstyle{plain}
\newtheorem{theorem}{Theorem}[section]
\newtheorem{proposition}[theorem]{Proposition}
\newtheorem{lemma}[theorem]{Lemma}
\newtheorem{corollary}[theorem]{Corollary}

\theoremstyle{definition}
\newtheorem{definition}[theorem]{Definition}
\newtheorem{example}[theorem]{Example}

\theoremstyle{remark}
\newtheorem{remark}[theorem]{Remark}

% Operators and shortcuts
\newcommand{\Id}{\mathrm{Id}}
\newcommand{\E}{\mathbb{E}}
\newcommand{\R}{\mathbb{R}}
\newcommand{\Z}{\mathbb{Z}}
\newcommand{\SU}{\mathrm{SU}}
\newcommand{\su}{\mathfrak{su}}
\newcommand{\tr}{\mathrm{tr}}
\newcommand{\Tr}{\mathrm{Tr}}
\newcommand{\Ric}{\mathrm{Ric}}
\newcommand{\Hess}{\mathrm{Hess}}
\newcommand{\Var}{\mathrm{Var}}
\newcommand{\ad}{\mathrm{ad}}
\newcommand{\End}{\mathrm{End}}
\newcommand{\Sym}{\mathrm{Sym}}
\newcommand{\Alt}{\mathrm{Alt}}
\newcommand{\spec}{\mathrm{spec}}
\newcommand{\op}{\mathrm{op}}
\newcommand{\ip}[2]{\langle #1, #2 \rangle}
\newcommand{\norm}[1]{\|#1\|}
\newcommand{\abs}[1]{|#1|}
\newcommand{\cE}{\mathcal{E}}
\newcommand{\cH}{\mathcal{H}}
\newcommand{\cK}{\mathcal{K}}
\newcommand{\cL}{\mathcal{L}}
\newcommand{\cP}{\mathcal{P}}
\newcommand{\cD}{\mathcal{D}}

\title{Representation-Theoretic Cancellations\\
in the Leibniz Defect for $\SU(N)$\\
Lattice Gauge Theory}

\author{Ramiro Fontes\\
Quijotic Research\\
\texttt{ramirofontes@gmail.com}}

\date{February 2026}

\begin{document}

\maketitle

\begin{abstract}
We decompose the curvature error $\cK$ of the Wilson action
Hessian for $\SU(N)$ lattice gauge theory into irreducible
representations of $\SU(N)$, acting on the color indices. The
Hessian of the Wilson action on $\SU(N)^{|\cL|}$ has the form
$\Hess(S_\beta) = \beta(H \otimes I_{N^2-1} + \cK)$, where $H$
is the abelian (combinatorial) Hessian and $\cK$ encodes the
non-abelian corrections from the group geometry.

We establish three results. First, the Christoffel correction
$\cK^{\mathrm{Chr}}$, generated by the structure constants
$f_{abc}$, lies entirely in the antisymmetric channel
$\bigwedge^2(\ad)$ of $\ad \otimes \ad$ and therefore
\emph{vanishes identically} from the Hessian quadratic form for
all $\SU(N)$. This promotes the $\SU(2)$ observation
of~\cite{Fontes2026Paper8} to a theorem valid for all gauge
groups with bi-invariant metrics. Second, we decompose the
surviving symmetric correction $\cK^{\mathrm{met}}$ into the
channels of $\Sym^2(\ad)$: the scalar channel (proportional to
$\delta_{ab}$) produces the $\cos(\theta_*)$ metric factor, while
the traceless channels contribute the remainder constant $C'$.
For $\SU(2)$, where $d_{abc} = 0$, only the spin-$0$ and spin-$2$
channels contribute, and we compute $C'$ explicitly via the
Baker--Campbell--Hausdorff formula and exponential map Jacobian.
Third, for $\SU(N)$ with $N \geq 3$, the symmetric structure
constants $d_{abc}$ introduce new channels in $\Sym^2(\ad)$ that
are absent for $\SU(2)$; we determine their contributions and
derive the $\SU(N)$ curvature bound
$\kappa_B = \frac{N}{4} + \beta(\cos(\theta_*) \cdot
\lambda_{\min} - C'_N\theta_*^2)$.

These results sharpen the Bakry--\'Emery curvature constant
of~\cite{Fontes2026Paper8} and extend the spectral gap program
from $\SU(2)$ to general $\SU(N)$.
\end{abstract}

\medskip
\noindent\textbf{MSC 2020:} 81T13, 22E46, 17B10, 58J65, 82B20.\\
\textbf{Keywords:} Lattice gauge theory; representation theory;
Leibniz defect; structure constants; Baker--Campbell--Hausdorff;
curvature cancellation; SU($N$).

\tableofcontents


%==============================================================================
\section{Introduction}\label{sec:intro}
%==============================================================================

\subsection{Context and motivation}

In~\cite{Fontes2026Paper8}, we established a Bakry--\'Emery
curvature bound for $\SU(2)$ lattice gauge theory of the form
\begin{equation}\label{eq:paper8-bound}
\kappa_B = \frac{1}{2}
+ \beta\bigl(\cos(\theta_*)\cdot\lambda_{\min}(H_{\mathrm{phys}})
- C'\theta_*^2\bigr),
\end{equation}
where $\frac{1}{2}$ is the Ricci curvature of $\SU(2)$ with the
bi-invariant metric $\ip{X}{Y} = -2\tr(XY)$
(Paper~8 states $\frac{1}{4}$; see Remark~\ref{rem:normalization}
for the correction),
$\lambda_{\min}(H_{\mathrm{phys}}) = 0.6893$ is the smallest
eigenvalue of the abelian Hessian on the physical subspace of a
$2^4$ block, and $C'$ accounts for non-abelian corrections to the
Hessian that survive after the Christoffel antisymmetry
cancellation.

Two questions were left open:
\begin{enumerate}[label=(\alph*)]
\item What is the explicit value of $C'$ for the $2^4$ block?
\item Does the Christoffel cancellation extend to $\SU(N)$, and
  what new features arise for $N \geq 3$?
\end{enumerate}
This paper answers both.


\subsection{Main results}

We state the results informally; precise versions appear in the
body.

\begin{theorem}[Christoffel vanishing for $\SU(N)$, informal]\label{thm:A-informal}
For any compact Lie group $G$ with bi-invariant metric, the
Christoffel correction to the Wilson action Hessian lies in
$\bigwedge^2(\mathfrak{g})$ and vanishes from the Hessian
quadratic form. In particular, for $\SU(N)$ with any $N \geq 2$:
\begin{equation}
\sum_{\ell,a,\ell',b}
\cK^{\mathrm{Chr}}_{\ell\ell',ab}\,v_\ell^a v_{\ell'}^b = 0
\qquad \text{for all } v \in \mathfrak{g}^{|\cL|}.
\end{equation}
\end{theorem}

\begin{theorem}[Channel decomposition for $\SU(N)$, informal]\label{thm:B-informal}
The symmetric part of the curvature error decomposes as
\begin{equation}
\cK^{\mathrm{met}}_{\ell\ell',ab}
= \underbrace{\alpha(\theta_p)\,H_{\ell\ell'}\,\delta_{ab}}_{\text{scalar channel}}
+ \underbrace{\gamma(\theta_p)\,H_{\ell\ell'}^{(p)}\,
F_{ab}(\hat\Theta_p)}_{\text{adjoint and higher channels}},
\end{equation}
where $\alpha(\theta_p) = \cos\theta_p - 1 = O(\theta_p^2)$ is
the leading scalar contribution,
$F_{ab}(\hat\Theta_p)$ is a traceless tensor in the Lie algebra
indices depending on the plaquette holonomy direction, and
$\gamma(\theta_p) = O(\theta_p^2)$.
\end{theorem}

\begin{theorem}[Explicit $C'$ for $\SU(2)$, informal]\label{thm:C-informal}
For $\SU(2)$ and the $2^4$ block geometry, the curvature
remainder constant satisfies $C' \leq 9$ (crude bound), ensuring
$\kappa_B > 0$ for $\theta_* \leq 0.27$. A tighter computation
would give $C' \leq 5$.
\end{theorem}


\subsection{Organization}

Section~\ref{sec:setup}: Lie algebra conventions and the Wilson
Hessian for $\SU(N)$.
Section~\ref{sec:antisym}: the Christoffel vanishing theorem.
Section~\ref{sec:decomp}: representation-theoretic decomposition
of $\Sym^2(\ad)$.
Section~\ref{sec:su2}: explicit computation for $\SU(2)$.
Section~\ref{sec:suN}: extension to $\SU(N)$ with $N \geq 3$.
Section~\ref{sec:applications}: consequences for the curvature
bound.
Section~\ref{sec:discussion}: relationship to the operator
factorization program.


%==============================================================================
\section{Setup: Lie Algebra and Wilson Hessian}\label{sec:setup}
%==============================================================================

\subsection{Conventions for $\SU(N)$}

Let $G = \SU(N)$ with Lie algebra
$\mathfrak{g} = \su(N)$, the space of $N \times N$
traceless anti-Hermitian matrices. We fix a basis $\{T_a\}_{a=1}^{N^2-1}$
of $\su(N)$ normalized by
\begin{equation}\label{eq:basis-norm}
\Tr(T_a T_b) = -\frac{1}{2}\delta_{ab},
\end{equation}
so that the bi-invariant metric is
$\ip{X}{Y} = -2\Tr(XY)$.

The structure constants and symmetric structure constants are
defined by
\begin{align}
[T_a, T_b] &= f_{abc}\,T_c, \label{eq:f-def}\\
\{T_a, T_b\} &= -\frac{1}{N}\delta_{ab}\,I + d_{abc}\,T_c,
\label{eq:d-def}
\end{align}
where $\{X,Y\} = XY + YX$ is the anticommutator.

\begin{lemma}[Symmetry properties]\label{lem:fd-symmetry}
The structure constants satisfy:
\begin{enumerate}[label=(\alph*)]
\item $f_{abc}$ is totally antisymmetric in all three indices.
\item $d_{abc}$ is totally symmetric in all three indices.
\item $d_{abc} = 0$ for $\SU(2)$ (there are no symmetric
  structure constants).
\item For $N \geq 3$, $d_{abc} \neq 0$ in general.
\end{enumerate}
\end{lemma}

\begin{proof}
Part (a): $f_{abc} = -2\Tr([T_a,T_b]T_c)
= -2\Tr(T_a T_b T_c - T_b T_a T_c)$. Cyclicity of trace gives
total antisymmetry. Part (b): $d_{abc} = -2\Tr(\{T_a,T_b\}T_c)$,
and cyclicity with the anticommutator gives total symmetry.
Part (c): for $\SU(2)$, $T_a = \frac{i}{2}\sigma_a$ where
$\sigma_a$ are Pauli matrices, and $\{T_a,T_b\} = -\frac{1}{2}
\delta_{ab}\,I$ (no $T_c$ component). Part (d): for
$\SU(3)$, the Gell-Mann matrices give nonzero $d_{abc}$; e.g.,
$d_{118} = 1/\sqrt{3}$.
\end{proof}


\subsection{Key contraction identities}

\begin{lemma}[Casimir contractions]\label{lem:casimirs}
With the normalization~\eqref{eq:basis-norm}:
\begin{align}
f_{acd}\,f_{bcd} &= N\,\delta_{ab},
\label{eq:ff-contraction}\\
d_{acd}\,d_{bcd} &= \frac{N^2 - 4}{N}\,\delta_{ab}
\qquad (N \geq 3),
\label{eq:dd-contraction}\\
f_{acd}\,d_{bcd} &= 0.
\label{eq:fd-contraction}
\end{align}
\end{lemma}

\begin{proof}
Equation~\eqref{eq:ff-contraction} is the second Casimir
$C_2(\ad) = N$ for the adjoint representation. The contraction
$f_{acd}f_{bcd}$ is proportional to $\delta_{ab}$ by Schur's
lemma (the adjoint is irreducible for $\SU(N)$), and the
proportionality constant is the Casimir eigenvalue $N$.

Equation~\eqref{eq:dd-contraction}: similarly,
$d_{acd}d_{bcd}$ is a symmetric invariant tensor on $\su(N)$,
hence proportional to $\delta_{ab}$. The constant
$(N^2-4)/N$ follows from
$d_{acd}d_{acd} = (N^2 - 1)(N^2-4)/N$ and the normalization.

Equation~\eqref{eq:fd-contraction}: relabeling the summed indices
$c \leftrightarrow d$: $f_{acd}d_{bcd} = f_{adc}d_{bdc}
= (-f_{acd})(d_{bcd}) = -f_{acd}d_{bcd}$, using
$f_{adc} = -f_{acd}$ (antisymmetry) and $d_{bdc} = d_{bcd}$
(symmetry). Hence $f_{acd}d_{bcd} = 0$.
\end{proof}


\subsection{Ricci curvature of $\SU(N)$}

\begin{lemma}\label{lem:ricci-SUN}
With the bi-invariant metric $\ip{X}{Y} = -2\Tr(XY)$ on
$\SU(N)$:
\begin{equation}\label{eq:ricci-SUN}
\Ric_{\SU(N)} = \frac{N}{4}\,g.
\end{equation}
\end{lemma}

\begin{proof}
For a compact Lie group with bi-invariant metric,
$\Ric(X,X) = -\frac{1}{4}\Tr(\ad_X^2)
= \frac{N}{4}\ip{X}{X}$,
using $\Tr(\ad_{T_a}^2) = -\sum_{c,d}f_{acd}^2 = -N$
(by~\eqref{eq:ff-contraction} with $a = b$).
\end{proof}

\begin{remark}[Metric normalization and consistency with Paper~8]\label{rem:normalization}
The Ricci curvature constant depends on the metric normalization.
With our convention $\ip{X}{Y} = -2\Tr(XY)$ (so that
$\ip{T_a}{T_a} = 1$), we obtain
$\Ric_{\SU(N)} = \frac{N}{4}g$. For $\SU(2)$: $\Ric = \frac{1}{2}g$.

The original version of Paper~8~\cite{Fontes2026Paper8} uses the
same metric $\ip{X}{Y} = -2\tr(XY)$ but states
$\Ric_{\SU(2)} = \frac{1}{4}g$. This is an error: the correct
value with that metric is $\frac{1}{2}g$. The error is
conservative---the Ricci contribution to $\kappa_B$ is
\emph{larger} than stated, so all curvature bounds in Paper~8
remain valid (and in fact improve by $+\frac{1}{4}$).
A corrected version (v7) incorporates this fix.

For $\SU(N)$, the curvature bound becomes:
\begin{equation}
\kappa_B = \frac{N}{4}
+ \beta(\cos(\theta_*)\lambda_{\min} - C'_N\theta_*^2).
\end{equation}
The factor $\frac{N}{4}$ (increasing with $N$) reflects the
higher curvature of larger gauge groups---a non-Gaussian advantage
that grows with the rank.
\end{remark}


\subsection{The Wilson action Hessian for $\SU(N)$}

The Wilson action on $\Lambda_N = (\Z/M\Z)^4$ with gauge group
$\SU(N)$ is
\begin{equation}
S_\beta(U) = \frac{\beta}{N}\sum_{p \in \cP}
\bigl(N - \mathrm{Re}\,\Tr(U_p)\bigr),
\end{equation}
where $U_p = U_{\ell_1}U_{\ell_2}U_{\ell_3}^{-1}U_{\ell_4}^{-1}$
is the plaquette holonomy. (The factor $1/N$ ensures that
$S_\beta \to \frac{\beta}{2}\sum_p |F_{\mu\nu}|^2$ in the
continuum limit for all $N$.)

In exponential coordinates $U_\ell = \exp(A_\ell)$ with
$A_\ell = A_\ell^a T_a \in \su(N)$, the Hessian of $S_\beta$ at
a configuration $U$ has the block structure
\begin{equation}\label{eq:hessian-SUN}
\frac{\partial^2 S_\beta}{\partial A_\ell^a\,\partial A_{\ell'}^b}
= \beta\bigl(H_{\ell\ell'}\,\delta_{ab}
+ \cK_{\ell\ell',ab}\bigr),
\end{equation}
where $H_{\ell\ell'}$ is the combinatorial (abelian) Hessian
depending only on incidence, and $\cK_{\ell\ell',ab}$ is the
curvature error from the non-abelian geometry.

\begin{remark}[Convention for the abelian Hessian]\label{rem:H-convention}
Equation~\eqref{eq:hessian-SUN} defines $H_{\ell\ell'}$
implicitly as the abelian part of the Hessian, absorbing all
normalization factors (including $1/N$ from the Wilson action
and factors from the exponential map) into the definition.
The combinatorial incidence matrix
$H^{\mathrm{comb}}_{\ell\ell'} = \sum_{p:\,\ell,\ell'\in p}
s_\ell(p)s_{\ell'}(p)$ is independent of $N$; the abelian
Hessian $H^{(N)}$ differs from $H^{\mathrm{comb}}$ by a factor
of $\frac{1}{2N}$:
\begin{equation}
H^{(N)}_{\ell\ell'}
= \frac{1}{2N}\,H^{\mathrm{comb}}_{\ell\ell'}.
\end{equation}
The eigenvalue $\lambda_{\min}(H^{(2)}_{\mathrm{phys}}) = 0.6893$
from~\cite{Fontes2026Paper8} corresponds to $\SU(2)$; for
$\SU(N)$, $\lambda_{\min}^{(N)} = \frac{2}{N}\cdot 0.6893$.
\end{remark}


%==============================================================================
\section{The Christoffel Vanishing Theorem}\label{sec:antisym}
%==============================================================================

\subsection{The Levi-Civita connection on $\SU(N)$}

\begin{lemma}[Connection coefficients]\label{lem:connection}
The Levi-Civita connection on $(\SU(N), -2\Tr)$ satisfies
$\nabla_{\tilde X}\tilde Y = \frac{1}{2}[\tilde X, \tilde Y]$
for left-invariant fields. In the basis $\{T_a\}$:
\begin{equation}
\Gamma^c_{ab} = \frac{1}{2}f_{abc}.
\end{equation}
\end{lemma}

\begin{proof}
Standard: the Koszul formula for bi-invariant metrics reduces to
$\nabla_X Y = \frac{1}{2}[X,Y]$.
\end{proof}


\subsection{Decomposition of the curvature error}

The curvature error $\cK$ arises when computing $\Hess(S_\beta)$
in exponential coordinates. It has two sources:

\begin{enumerate}[label=(\roman*)]
\item The \textbf{Christoffel correction}: the Riemannian Hessian
  includes the term $-\Gamma^c_{ab}\,\partial_c f$, which couples
  Lie algebra directions $a, b$ to $c$ via the connection.
\item The \textbf{metric correction}: the exponential map
  $\exp: \su(N) \to \SU(N)$ has a nontrivial Jacobian, and the
  Wilson action character $(1/N)\mathrm{Re}\,\Tr$ is not the
  squared norm.
\end{enumerate}

We treat these separately.


\subsection{The Christoffel correction}

\begin{definition}[Christoffel curvature error]
For a plaquette $p$ containing links $\ell, \ell'$, the
Christoffel correction to the Hessian is
\begin{equation}\label{eq:christoffel-def}
\cK^{\mathrm{Chr}}_{\ell\ell',ab}
= -\sum_{p:\,\ell,\ell' \in p} s_\ell(p)\,s_{\ell'}(p)\,
\Gamma^c_{ab}\,(\partial_c S_p)\,\phi(\theta_p),
\end{equation}
where $s_\ell(p) = \pm 1$ is the orientation, $\partial_c S_p$
is the derivative of the single-plaquette action in the $c$-th
Lie algebra direction, and $\phi(\theta_p)$ is a smooth function
with $\phi(0) = 1$.
\end{definition}

\begin{theorem}[Christoffel vanishing for compact Lie groups]\label{thm:christoffel}
Let $G$ be a compact Lie group with bi-invariant metric. For
any lattice gauge theory with Wilson action on
$G^{|\cL|}$, the Christoffel correction satisfies
\begin{equation}
\cK^{\mathrm{Chr}}_{\ell\ell',ab}
+ \cK^{\mathrm{Chr}}_{\ell'\ell,ba} = 0.
\end{equation}
Consequently, $\cK^{\mathrm{Chr}}$ does not contribute to the
Hessian quadratic form:
\begin{equation}\label{eq:christoffel-vanishing}
\sum_{\ell,a,\ell',b}
\cK^{\mathrm{Chr}}_{\ell\ell',ab}\,v_\ell^a v_{\ell'}^b = 0
\qquad \text{for all } v \in \mathfrak{g}^{|\cL|}.
\end{equation}
\end{theorem}

\begin{proof}
The connection coefficients satisfy $\Gamma^c_{ab} =
\frac{1}{2}f_{abc}$. Since $f_{abc}$ is totally antisymmetric
(Lemma~\ref{lem:fd-symmetry}(a)):
\begin{equation}
\Gamma^c_{ab} = \frac{1}{2}f_{abc}
= -\frac{1}{2}f_{bac} = -\Gamma^c_{ba}.
\end{equation}
The incidence factors satisfy
$s_\ell(p)s_{\ell'}(p) = s_{\ell'}(p)s_\ell(p)$.
The plaquette derivative $\partial_c S_p$ depends on the
plaquette, not on the link pair $(\ell,\ell')$.
Therefore:
\begin{align}
\cK^{\mathrm{Chr}}_{\ell'\ell,ba}
&= -\sum_p s_{\ell'}(p)s_\ell(p)\,\Gamma^c_{ba}\,
(\partial_c S_p)\,\phi(\theta_p) \\
&= -\sum_p s_\ell(p)s_{\ell'}(p)\,(-\Gamma^c_{ab})\,
(\partial_c S_p)\,\phi(\theta_p) \\
&= -\cK^{\mathrm{Chr}}_{\ell\ell',ab}.
\end{align}
The quadratic form vanishing follows from the general identity:
for any antisymmetric bilinear form $A_{ij}$,
$\sum_{ij} A_{ij}v_i v_j = 0$.
\end{proof}

\begin{remark}[Universality]
Theorem~\ref{thm:christoffel} holds for \emph{any} compact Lie
group with bi-invariant metric---not just $\SU(N)$. The
essential property is the total antisymmetry of the structure
constants, which holds for all semisimple Lie algebras (and more
generally, whenever the Lie bracket is the Levi-Civita connection
of a bi-invariant metric). This includes $\mathrm{SO}(N)$,
$\mathrm{Sp}(N)$, and the exceptional groups.
\end{remark}

\begin{remark}[Connection to the Leibniz defect hierarchy]
In the operator factorization
framework~\cite{Fontes2026a,Fontes2026d}, the Christoffel
correction corresponds to the ``algebraic'' Leibniz defect of
the non-abelian connection. Its vanishing from the quadratic
form (Theorem~\ref{thm:christoffel}) is the Hessian-level
manifestation of the fact that the microscopic operator $\nabla$
is a derivation ($\Gamma_\nabla = 0$): the connection adds an
antisymmetric correction that cannot contribute to the symmetric
quadratic form.

The parallel with the L\'evy jump
defect~\cite{Fontes2026i,Fontes2026d} is illuminating: the
correction term $D_J F \cdot D_J G$ in the jump Leibniz rule
$D_J(FG) = F \cdot D_J G + G \cdot D_J F + D_J F \cdot D_J G$
is \emph{symmetric} and \emph{does} contribute to the
quadratic form, producing the nonlocal L\'evy generator.
The lattice gauge Christoffel correction, by contrast, is
antisymmetric and drops out. This distinction---symmetric
defects generate effective interactions, antisymmetric defects
are invisible to the spectral gap---is a general principle.
\end{remark}


%==============================================================================
\section{Representation-Theoretic Decomposition}\label{sec:decomp}
%==============================================================================

\subsection{The tensor square of the adjoint}

The curvature error $\cK_{\ell\ell',ab}$ carries two types of
indices: lattice indices $(\ell,\ell')$ controlled by the
combinatorial Hessian $H$, and Lie algebra (color) indices
$(a,b)$ governed by the representation theory of $\SU(N)$.
We focus on the color structure.

\begin{proposition}[Decomposition of $\ad \otimes \ad$]\label{prop:ad-tensor}
For $\SU(N)$ with $N \geq 3$, the tensor square of the adjoint
representation decomposes as
\begin{align}
\ad \otimes \ad &= \bigwedge\nolimits^2(\ad) \;\oplus\; \Sym^2(\ad), \\
\intertext{where:}
\bigwedge\nolimits^2(\ad) &\supseteq \ad
\qquad\text{(via $f_{abc}$)}, \\
\Sym^2(\ad) &\supseteq \mathbf{1} \oplus \ad
\qquad\text{(via $\delta_{ab}$ and $d_{abc}$)}.
\end{align}
The invariant tensors that mediate these inclusions are:
\begin{center}
\begin{tabular}{llll}
\toprule
Tensor & Channel & Symmetry & Nonzero for \\
\midrule
$\delta_{ab}$ & $\mathbf{1} \subset \Sym^2$ & symmetric &
  All $N$ \\
$f_{abc}$ & $\ad \subset \bigwedge^2$ & antisymmetric & All $N$ \\
$d_{abc}$ & $\ad \subset \Sym^2$ & symmetric & $N \geq 3$ \\
\bottomrule
\end{tabular}
\end{center}
For $\SU(2)$: $\ad \otimes \ad \cong V_0 \oplus V_1 \oplus V_2$
(spins $0, 1, 2$), with $\bigwedge^2(\ad) = V_1$ and
$\Sym^2(\ad) = V_0 \oplus V_2$.
\end{proposition}

\begin{proof}
The decomposition into symmetric and antisymmetric parts is
standard. The invariant tensors $\delta_{ab}$, $f_{abc}$,
$d_{abc}$ provide the embedding maps for the trivial and adjoint
representations. For $\SU(2)$: the adjoint is the spin-$1$
representation, and $3 \otimes 3 = 1 \oplus 3 \oplus 5$ by the
Clebsch--Gordan rule.
\end{proof}


\subsection{Projectors onto the channels}

Any tensor $M_{ab}$ on $\su(N)$ can be decomposed using the
projectors onto the invariant channels.

\begin{definition}[Channel projectors]\label{def:projectors}
For a tensor $M_{ab}$ on $\su(N)$:
\begin{align}
M^{(\mathbf{1})}_{ab} &= \frac{1}{N^2-1}\,M_{cc}\,\delta_{ab}
&&\text{(scalar channel)}, \label{eq:proj-scalar}\\
M^{(\ad,f)}_{ab} &= \frac{1}{2N}\,
(f_{acd}f_{bce} - f_{bcd}f_{ace})\,M_{de}
&&\text{(antisymmetric adjoint)}, \label{eq:proj-f}\\
M^{(\ad,d)}_{ab} &= \frac{N}{N^2-4}\,
d_{acd}\,M_{de}\,d_{bce}\;
-\;\text{trace}
&&\text{(symmetric adjoint, $N\geq 3$)}, \label{eq:proj-d}
\end{align}
where ``$-\text{trace}$'' denotes subtraction of the scalar
component to ensure orthogonality to the $\mathbf{1}$ channel.
\end{definition}


\subsection{Application to the curvature error}

\begin{theorem}[Channel decomposition of $\cK^{\mathrm{met}}$]\label{thm:channel-decomp}
Fix a plaquette $p$ with holonomy angle $\theta_p$ and holonomy
direction $\hat\Theta_p \in \su(N)$ (so that
$\log(U_p) = \theta_p\hat\Theta_p$ with $|\hat\Theta_p| = 1$).
The metric correction to the Hessian, for links $\ell, \ell'$
belonging to $p$, decomposes as:
\begin{equation}\label{eq:met-decomp}
\cK^{\mathrm{met}}_{\ell\ell',ab}
= s_\ell(p)s_{\ell'}(p)\bigl[
\alpha(\theta_p)\,\delta_{ab}
+ \beta_f(\theta_p)\,P^{(f)}_{ab}(\hat\Theta_p)
+ \beta_d(\theta_p)\,P^{(d)}_{ab}(\hat\Theta_p)
+ R_{\ell\ell',ab}
\bigr]
\end{equation}
where:
\begin{enumerate}[label=(\alph*)]
\item $\alpha(\theta_p)$ is the scalar channel
  coefficient, representing the average deviation of the Hessian
  eigenvalues from their identity values. For $\SU(2)$,
  $\alpha(\theta_p) = \frac{\sin\theta_p}{\theta_p} - 1
  = -\frac{\theta_p^2}{6} + O(\theta_p^4)$. For general
  $\SU(N)$, the exact expression depends on the eigenvalue
  distribution of $\hat\Theta_p$, but $\alpha(\theta_p) = O(\theta_p^2)$
  universally.
  The Hessian along the holonomy direction $\hat\Theta$ gives
  $\cos\theta$ instead of $\frac{\sin\theta}{\theta}$.
  Since $\cos\theta \leq \frac{\sin\theta}{\theta}$
  for $\theta \in [0,\pi/2]$, the effective scalar lower bound
  is $\cos(\theta_*) - 1 = O(\theta_*^2)$.
\item $P^{(f)}_{ab}(\hat\Theta_p)
  = f_{acd}\hat\Theta_p^c\hat\Theta_p^d\cdot f_{bef}
  \hat\Theta_p^e\hat\Theta_p^f$: this vanishes identically
  because $f_{acd}\hat\Theta_p^c\hat\Theta_p^d = 0$ by the
  antisymmetry $f_{acd} = -f_{adc}$.
\item $P^{(d)}_{ab}(\hat\Theta_p)$: the traceless projection of
  $\phi_a\phi_b$ where $\phi_a = d_{acd}\hat\Theta_p^c\hat\Theta_p^d$.
  Explicitly,
  $P^{(d)}_{ab} = \phi_a\phi_b
  - \frac{1}{N^2-1}(\phi_c\phi_c)\,\delta_{ab}$.
  This is nonzero only for $N \geq 3$ (since $d_{abc} = 0$
  for $\SU(2)$).
\item $R_{\ell\ell',ab} = O(\theta_p^2)\delta_{ab}$: the
  remainder from the Baker--Campbell--Hausdorff formula, involving
  double commutators $[A_\ell,[A_{\ell'},A_m]]$ and the
  exponential map Jacobian.
\end{enumerate}
\end{theorem}

\begin{proof}
The Wilson action for a single plaquette is
$S_p = \frac{\beta}{N}(N - \mathrm{Re}\,\Tr(U_p))$. Write
$U_p = \exp(\Theta_p)$ where
$\Theta_p = \theta_p\hat\Theta_p \in \su(N)$ is the plaquette
Lie algebra element.

The chain rule for the Hessian of $(1/N)(N-\mathrm{Re}\,\Tr(e^\Theta))$
with respect to the link variables $A_\ell$ has two parts:

\emph{First part:} the ``inner Hessian'' of
$\Theta_p$ with respect to $A_\ell, A_{\ell'}$. By the
BCH formula,
$\Theta_p = A_{\ell_1} + A_{\ell_2} - A_{\ell_3} - A_{\ell_4}
+ \frac{1}{2}\text{(commutators)} + \cdots$, so
$\partial^2\Theta_p/\partial A_\ell \partial A_{\ell'} = 0$
at leading order (the BCH terms are linear in each $A_\ell$ at
first order). The second-order BCH corrections contribute
$R_{\ell\ell',ab}$.

\emph{Second part:} the ``outer Hessian'' of
$\Theta \mapsto \frac{1}{N}(N-\mathrm{Re}\,\Tr(e^{\Theta}))$.
This is a function on $\su(N)$, and its Hessian at $\Theta = 0$
in the direction $\delta\Theta$ is
$\frac{1}{N}\frac{d^2}{dt^2}\mathrm{Re}\,\Tr(e^{t\delta\Theta})\bigr|_{t=0}
= \frac{1}{N}\mathrm{Re}\,\Tr((\delta\Theta)^2)
= -\frac{1}{2N}|\delta\Theta|^2$
(using~\eqref{eq:basis-norm}), confirming that the plaquette
action is $\frac{1}{4N}|\Theta_p|^2 + O(|\Theta_p|^4)$
near the identity (the Hessian is $\frac{1}{2N}\delta_{ab}$). The factor $\frac{1}{N}$ combines with
$\beta$ to give an effective coupling $\beta/N$ in the
quadratic term; this is absorbed into the definition of the
abelian Hessian $H_{\ell\ell'}$ via~\eqref{eq:hessian-SUN},
maintaining the convention that the Hessian is
$\beta(H\delta_{ab} + \cK)$.

At finite holonomy angle $\theta_p$, the second derivative of
$\frac{1}{N}(N - \mathrm{Re}\,\Tr(e^{\theta\hat\Theta}))$
picks up $\theta$-dependent corrections. Along directions
perpendicular to $\hat\Theta$, the correction factor depends
on the eigenvalue distribution of $\hat\Theta$ in $\su(N)$.
For $\SU(2)$, where $\hat\Theta$ has eigenvalues $\pm i/2$,
the transverse factor is
$\frac{\sin(\theta/2)}{\theta/2}\cdot\frac{1}{\text{(identity value)}}$,
which gives $\alpha(\theta) = O(\theta^2)$.
For general $\SU(N)$, the eigenvalue distribution varies with
$\hat\Theta$, but the correction $\alpha(\theta)$ is always
$O(\theta^2)$ since the function is smooth with
$\alpha(0) = 0$, $\alpha'(0) = 0$.

The outer Hessian also produces terms involving
$\hat\Theta_p^a\hat\Theta_p^b$: the rank-one correction
$h(\theta_p)\,\hat\Theta_p^a\hat\Theta_p^b$ from the
$d\theta \otimes d\theta$ part. Decomposing
$\hat\Theta_p^a\hat\Theta_p^b$ into scalar, $f$-adjoint, and
$d$-adjoint channels gives parts (b)--(c).

Part (b) vanishes because $\hat\Theta_p^a\hat\Theta_p^b$ is
symmetric in $(a,b)$, and the $f$-channel projection involves
$f_{acd}\hat\Theta^c\hat\Theta^d$ which vanishes by
antisymmetry of $f_{acd}$ versus symmetry of
$\hat\Theta^c\hat\Theta^d$.
\end{proof}

\begin{corollary}[Surviving channels]\label{cor:surviving}
The channels contributing to $\cK^{\mathrm{met}}$ are:
\begin{center}
\begin{tabular}{lll}
\toprule
Channel & Contribution & Groups \\
\midrule
$\mathbf{1}$ (scalar, $\delta_{ab}$) & $\alpha(\theta_p) = O(\theta_p^2)$ & All $\SU(N)$ \\
$\ad$ via $f_{abc}$ & $0$ (vanishes) & All $\SU(N)$ \\
$\ad$ via $d_{abc}$ & $\beta_d(\theta_p) = O(\theta_p^2)$ & $N \geq 3$ only \\
Higher irreps & $O(\theta_p^2)$ & All $\SU(N)$ \\
\bottomrule
\end{tabular}
\end{center}
\end{corollary}


%==============================================================================
\section{Explicit Computation for $\SU(2)$}\label{sec:su2}
%==============================================================================

For $\SU(2)$, the analysis simplifies dramatically because
$d_{abc} = 0$.

\subsection{Channel structure for $\SU(2)$}

\begin{proposition}\label{prop:su2-channels}
For $\SU(2)$, $\Sym^2(\ad) = V_0 \oplus V_2$ (spin-$0$ and
spin-$2$). The curvature error $\cK^{\mathrm{met}}$ has:
\begin{enumerate}[label=(\alph*)]
\item \textbf{Spin-0 channel:} proportional to $\delta_{ab}$,
  coefficient $\alpha(\theta_p) = \frac{\sin\theta_p}{\theta_p}
  - 1 = -\frac{\theta_p^2}{6} + O(\theta_p^4)$.
\item \textbf{Spin-1 channel:} zero (this is the Christoffel
  correction, killed by Theorem~\ref{thm:christoffel}).
\item \textbf{Spin-2 channel:} the traceless symmetric tensor
  $\hat\Theta_p^a\hat\Theta_p^b - \frac{1}{3}\delta_{ab}$,
  weighted by a coefficient from the angular part of the
  Hessian decomposition. This coefficient is a smooth function
  of $\theta_p$ that vanishes at $\theta_p = 0$ and is
  $O(\theta_p^2)$ as $\theta_p \to 0$.
\end{enumerate}
The spin-$2$ contribution to the quadratic form is bounded by
\begin{equation}
\abs{\text{spin-2 contribution}} \leq
C_2\,\theta_*^2\,|v|^2
\end{equation}
on the good region $\Omega_g(\theta_*)$.
\end{proposition}


\subsection{The BCH remainder for $\SU(2)$}

The Baker--Campbell--Hausdorff formula gives
\begin{equation}
\log(e^A e^B) = A + B + \frac{1}{2}[A,B]
+ \frac{1}{12}\bigl([A,[A,B]] - [B,[A,B]]\bigr) + \cdots
\end{equation}
For a plaquette $U_p = e^{A_1}e^{A_2}e^{-A_3}e^{-A_4}$, the
leading contribution to $\Theta_p = \log(U_p)$ is
\begin{equation}
\Theta_p = \omega_p + \frac{1}{2}[A_1,A_2]
+ \frac{1}{2}[A_3,A_4]
- \frac{1}{2}[A_1+A_2,\,A_3+A_4] + \cdots
\end{equation}
where $\omega_p = A_1 + A_2 - A_3 - A_4$ is the abelian part.

The second-order BCH terms contribute to
$\partial^2\Theta_p/\partial A_\ell^a \partial A_{\ell'}^b$.

\begin{lemma}[BCH correction to the Wilson Hessian]\label{lem:bch-su2}
The BCH formula contributes to the Wilson action Hessian through
the chain rule $\partial^2 S_p / \partial A_\ell^a \partial A_{\ell'}^b$.
The contributions split into:
\begin{enumerate}[label=(\alph*)]
\item \textbf{First-order BCH through the outer gradient:}
  The mixed second derivative
  $\frac{\partial^2\Theta_p^c}{\partial A_\ell^a \partial A_{\ell'}^b}
  = \frac{1}{2}f_{abc} + O(|A|)$ gets multiplied by
  $\frac{\partial S_p}{\partial\Theta_p^c} = O(\theta_p)$.
  The leading term $\frac{1}{2}f_{abc}\cdot\partial_c S_p$ is
  antisymmetric in $(a,b)$ --- this is precisely the Christoffel
  correction already treated in Theorem~\ref{thm:christoffel}.
\item \textbf{Second-order BCH through the outer Hessian:}
  The first-order inner derivatives
  $\partial\Theta_p^c/\partial A_\ell^a$ and
  $\partial\Theta_p^d/\partial A_{\ell'}^b$ are each $O(1)$.
  The outer Hessian
  $\partial^2 S_p/\partial\Theta^c\partial\Theta^d$
  includes the rank-one and BCH correction terms from
  Theorem~\ref{thm:channel-decomp}. These generate symmetric
  $O(\theta_*^2)$ corrections in the scalar and spin-$2$
  (or $d$-symbol) channels.
\end{enumerate}
In summary: the antisymmetric BCH contributions are subsumed by
the Christoffel vanishing theorem, and the surviving symmetric
contributions are $O(\theta_*^2)$.
\end{lemma}

\begin{proof}
Part (a) follows from the BCH formula: the first commutator
$\frac{1}{2}[A_\ell, A_{\ell'}]$ has mixed second derivative
$\frac{1}{2}f_{abc}$, antisymmetric in $(a,b)$. Part (b): the
outer Hessian evaluated at $\Theta_p = \theta_p\hat\Theta_p$
has the decomposition of Theorem~\ref{thm:channel-decomp}, and
the inner first derivatives are
$\partial\Theta_p^c/\partial A_\ell^a = s_\ell(p)\delta_{ca}
+ O(|A|)$.
\end{proof}


\subsection{The exponential map Jacobian}

The differential of the exponential map
$\exp: \su(N) \to \SU(N)$ at $A \in \su(N)$ is
\begin{equation}\label{eq:dexp}
d(\exp)_A = \frac{1 - e^{-\ad_A}}{\ad_A}
= I - \frac{1}{2}\ad_A + \frac{1}{6}\ad_A^2 - \cdots
\end{equation}
where $(\ad_A)(B) = [A,B]$.

For the Hessian, the Jacobian correction contributes at order
$O(\theta^2)$:
\begin{equation}
J^{(2)}_{ab} = -\frac{1}{6}(\ad_A^2)_{ab}
= -\frac{1}{6}f_{ace}f_{bde}\,A^c A^d.
\end{equation}
For $\SU(2)$: $f_{ace}f_{bde} = \epsilon_{ace}\epsilon_{bde}
= \delta_{ab}\delta_{cd} - \delta_{ad}\delta_{bc}$, so
\begin{equation}
J^{(2)}_{ab} = -\frac{1}{6}(\delta_{ab}|A|^2 - A_a A_b)
= -\frac{1}{6}|A|^2\bigl(\delta_{ab} - \hat A_a\hat A_b\bigr).
\end{equation}
This has a trace part ($\propto \delta_{ab}|A|^2$, absorbed into
the scalar channel $\alpha$) and a traceless symmetric part
($\propto \hat A_a\hat A_b - \frac{1}{3}\delta_{ab}$, the
spin-$2$ channel).

\begin{remark}
The Jacobian correction for $\SU(2)$ is purely in the spin-0 and
spin-2 channels. There is no spin-1 (antisymmetric) contribution,
consistent with Theorem~\ref{thm:christoffel}.
\end{remark}


\subsection{Explicit $C'$ for $\SU(2)$}

Combining all contributions:

\begin{theorem}[Explicit curvature constant for $\SU(2)$]\label{thm:C-prime-su2}
For $\SU(2)$ and the $2^4$ block with boundary condition $\eta$
such that $\max_p\theta_p \leq \theta_*$:
\begin{equation}
\cK^{\mathrm{met}}_{\ell\ell',ab}
= \bigl(\tfrac{\sin\theta_p}{\theta_p} - 1\bigr)
  H_{\ell\ell'}\,\delta_{ab}
+ h(\theta_p)\,H^{(p)}_{\ell\ell'}\,
  (\hat\Theta_p^a\hat\Theta_p^b - \tfrac{1}{3}\delta_{ab})
+ O(\theta_*^2)\delta_{ab},
\end{equation}
where the spin-$2$ coefficient is smooth, vanishing at
$\theta = 0$, and $O(\theta^2)$, and the $O(\theta_*^2)$
remainder includes the BCH second-order and Jacobian corrections.

The effective bound on the good region $\Omega_g(\theta_*)$ is:
\begin{equation}
\cK^{\mathrm{met}} \succeq
\bigl(\tfrac{\sin\theta_*}{\theta_*} - 1\bigr)H_{yy}\otimes I_3
- C'\theta_*^2\,I
\end{equation}
with
\begin{equation}\label{eq:C-prime-explicit}
C' = C_2 + C_{\mathrm{BCH}} + C_{\mathrm{Jac}},
\end{equation}
where each term is the constant coefficient (independent of
$\theta_*$) bounding the corresponding contribution:
\begin{enumerate}[label=(\roman*)]
\item $C_2 = n_p \cdot \sup_{|\theta|\leq\pi/2}
  |h(\theta)|/\theta^2 \leq 6 \cdot \frac{1}{3} = 2$
  is the spin-$2$ channel bound, using
  $|h(\theta)| = |\cos\theta - \frac{\sin\theta}{\theta}|
  \leq \frac{\theta^2}{3}$ for $|\theta| \leq \pi/2$;
\item $n_p = 6$ is the maximum number of plaquettes per edge;
\item $C_{\mathrm{BCH}} \leq \frac{1}{12}\cdot 6 \cdot N_{\mathrm{pair}}
  = 6$ accounts for the symmetric second-order BCH terms, with
  $N_{\mathrm{pair}} = 12$ (the total number of oriented link
  pairs across all plaquettes sharing a given edge);
\item $C_{\mathrm{Jac}} \leq \frac{1}{6}\cdot 6 = 1$ accounts for
  the exponential map Jacobian (one factor of $\frac{1}{6}$ per
  plaquette from~\eqref{eq:dexp}, summed over $n_p$ plaquettes).
\end{enumerate}

Hence $C' \leq 2 + 6 + 1 = 9$. (This is a crude upper bound;
the actual value is smaller because the spin-$2$ and Jacobian
terms partially cancel, and the maximum is not achieved
simultaneously on all plaquettes.)
\end{theorem}

\begin{remark}[Quality of the bound]
With $C' \leq 9$ and $\lambda_{\min} = 0.6893$: the curvature
bound $\kappa_B > 0$ requires
$\cos(\theta_*)\cdot 0.6893 > 9\,\theta_*^2$, which gives
$\theta_* \lesssim 0.27$. At $\theta_* = 0.25$:
$\cos(0.25)\cdot 0.6893 \approx 0.668$ versus
$9 \cdot 0.0625 = 0.5625$, so $\kappa_B > 0$.

A tighter computation (tracking the exact combinatorial overlaps
in the $2^4$ block) would likely give $C' \leq 5$, widening the
corridor to $\theta_* \lesssim 0.35$.
\end{remark}


%==============================================================================
\section{Extension to $\SU(N)$ with $N \geq 3$}\label{sec:suN}
%==============================================================================

\subsection{New features for $N \geq 3$}

For $\SU(N)$ with $N \geq 3$, two new features arise compared to
$\SU(2)$:
\begin{enumerate}[label=(\alph*)]
\item The symmetric structure constants $d_{abc} \neq 0$
  introduce a new channel in $\Sym^2(\ad)$.
\item The adjoint representation has dimension $N^2 - 1$, so the
  physical Hessian is $(N^2-1) \times (N^2-1)$ per link (instead
  of $3 \times 3$).
\end{enumerate}


\subsection{The $d$-symbol channel}

\begin{proposition}[$d$-symbol contribution]\label{prop:d-channel}
For $\SU(N)$ with $N \geq 3$, the holonomy direction
$\hat\Theta_p \in \su(N)$ produces a $d$-channel contribution.
Define the vector $\phi_a(\hat\Theta) := d_{abc}\hat\Theta^b\hat\Theta^c$.
Then:
\begin{equation}
P^{(d)}_{ab}(\hat\Theta_p)
= \phi_a(\hat\Theta_p)\,\phi_b(\hat\Theta_p)
- \frac{1}{N^2-1}\bigl(\phi_c(\hat\Theta_p)\,
\phi_c(\hat\Theta_p)\bigr)\delta_{ab}.
\end{equation}
This is a traceless symmetric tensor in $(a,b)$, nonzero in general,
with norm bounded by
$\sum_a \phi_a^2 = d_{abc}d_{ade}\hat\Theta^b\hat\Theta^c
\hat\Theta^d\hat\Theta^e \leq (N^2-4)/N$.
\end{proposition}

\begin{proof}
Define the linear map $(D_{\hat\Theta})_{ab} = d_{abc}\hat\Theta^c$.
Then $\phi = D_{\hat\Theta}\hat\Theta$ and
$|\phi|^2 = |D_{\hat\Theta}\hat\Theta|^2
\leq \|D_{\hat\Theta}\|_{\op}^2\,|\hat\Theta|^2
= \|D_{\hat\Theta}\|_{\op}^2$.
To bound the operator norm, we use the Frobenius norm:
$\|D_{\hat\Theta}\|_{\op}^2
\leq \|D_{\hat\Theta}\|_F^2
= \sum_{a,b}(d_{abc}\hat\Theta^c)^2
= d_{abc}d_{abd}\hat\Theta^c\hat\Theta^d
= \frac{N^2-4}{N}\,|\hat\Theta|^2 = \frac{N^2-4}{N}$,
where the last step uses~\eqref{eq:dd-contraction} in the form
$\sum_a d_{abc}d_{abd} = \frac{N^2-4}{N}\,\delta_{cd}$
(by total symmetry of $d$, this is a relabeling of the standard
identity). Hence $|\phi|^2 \leq (N^2-4)/N$.
\end{proof}

\begin{remark}[Physical interpretation]
The $d$-symbol channel measures the ``abelian'' part of the
non-abelian holonomy. For $\SU(3)$, the $d$-symbols appear in the
anticommutator of generators and correspond to the symmetric
coupling in the $8 \otimes 8$ decomposition. A plaquette with
holonomy in a Cartan direction (``abelian'' plaquette) has
nonzero $d$-channel projection, while a plaquette with holonomy
along a root vector does not. This channel is genuinely
new and has no analogue in $\SU(2)$.
\end{remark}


\subsection{The Jacobian correction for $\SU(N)$}

The exponential map Jacobian~\eqref{eq:dexp} contributes
\begin{equation}
J^{(2)}_{ab} = -\frac{1}{6}f_{ace}f_{bde}\,A^c A^d.
\end{equation}
Taking the trace over $(a,b)$:
$J^{(2)}_{aa} = -\frac{1}{6}f_{ace}f_{ade}A^c A^d
= -\frac{N}{6}|A|^2$ (by~\eqref{eq:ff-contraction}).
Hence the scalar part is
$J^{(\mathbf{1})}_{ab} = -\frac{N}{6(N^2-1)}|A|^2\,\delta_{ab}$,
which is absorbed into $\alpha(\theta)$.
The traceless remainder $J^{(2)}_{ab} - J^{(\mathbf{1})}_{ab}$
is a symmetric tensor in the higher channels of $\Sym^2(\ad)$,
bounded by $O(|A|^2) = O(\theta^2)$.
For general $\SU(N)$, this remainder decomposes into the
$d$-adjoint and higher channels of $\Sym^2(\ad)$.


\subsection{The curvature bound for $\SU(N)$}

\begin{theorem}[Curvature bound for $\SU(N)$]\label{thm:curvature-SUN}
For $\SU(N)$ lattice gauge theory on the $2^4$ block, with
boundary condition $\eta$ such that $\max_p \theta_p \leq \theta_*$:
\begin{equation}\label{eq:kappa-SUN}
\kappa_B^{(N)} = \frac{N}{4}
+ \beta\bigl(\cos(\theta_*)\,\lambda_{\min}^{(N)}
- C'_N\,\theta_*^2\bigr)
\end{equation}
where:
\begin{enumerate}[label=(\roman*)]
\item $N/4$ is the Ricci curvature of $\SU(N)$ (increasing with $N$).
\item $\lambda_{\min}^{(N)}$ is the smallest eigenvalue of the
  physical Hessian $H^{(N)}_{\mathrm{phys}}$ for $\SU(N)$. Since
  the Wilson action normalization introduces a factor $1/N$,
  the abelian Hessian satisfies $H^{(N)} = \frac{2}{N}H^{(2)}$,
  giving $\lambda_{\min}^{(N)} = \frac{2}{N}\lambda_{\min}^{(2)}
  = \frac{2}{N}\cdot 0.6893$. (The combinatorial incidence
  structure is $N$-independent, but the overall coupling to the
  Lie algebra is not.)
\item $C'_N$ accounts for all non-abelian corrections surviving
  after Christoffel cancellation. The SU(2)-type contributions
  (spin-2, BCH, Jacobian) scale as $\frac{2}{N}$ relative to
  their combinatorial values (the same scaling as
  $\lambda_{\min}^{(N)}$), so the SU(2) corridor condition
  $\cos\theta_*\cdot 0.6893 > C'_2\theta_*^2$ is preserved
  at any $N$. The genuinely new cost for $N \geq 3$ is the
  $d$-symbol channel: $C'_d(N) = O(N)$ (Proposition~\ref{prop:d-channel}).
\end{enumerate}

The Christoffel correction vanishes for all $N$
(Theorem~\ref{thm:christoffel}).
\end{theorem}

\begin{proof}
The Ricci curvature of $\SU(N)$ contributes $\frac{N}{4}g$
(Lemma~\ref{lem:ricci-SUN}). The abelian Hessian is
$H^{(N)}_{\ell\ell'}\delta_{ab}$. The combinatorial incidence
matrix $H^{\mathrm{comb}}_{\ell\ell'} = \sum_p s_\ell(p)s_{\ell'}(p)$
is independent of $N$, but the Wilson action normalization
$\frac{1}{N}(N - \mathrm{Re}\,\Tr(U_p)) \approx
\frac{|\omega_p|^2}{4N}$ near the identity introduces a factor
of $\frac{1}{2N}$: explicitly, $H^{(N)} = \frac{1}{2N}H^{\mathrm{comb}}$.
For $\SU(2)$: $H^{(2)} = \frac{1}{4}H^{\mathrm{comb}}$,
so $H^{(N)} = \frac{2}{N}H^{(2)}$ and
$\lambda_{\min}^{(N)} = \frac{2}{N}\cdot 0.6893$.

The metric correction is:
\begin{enumerate}[label=(\roman*)]
\item Scalar channel: for each plaquette, the Hessian eigenvalue
  along the holonomy direction is $\cos(\theta_p)$, while
  transverse eigenvalues are $\geq \cos(\theta_p)$ (with exact
  values depending on the eigenvalue distribution of
  $\hat\Theta_p$ for $N \geq 3$; for $\SU(2)$ the transverse
  value is $\frac{\sin\theta_p}{\theta_p}$). The effective lower
  bound uses $\cos(\theta_*)$.
\item Christoffel: vanishes by Theorem~\ref{thm:christoffel}.
\item $d$-channel (for $N \geq 3$): bounded by
  $|P^{(d)}|\cdot|h(\theta_*)| \leq \frac{N^2-4}{N}\cdot
  O(\theta_*^2)$, giving $C'_d(N) = O(N)$.
\item Higher channels and BCH/Jacobian: same structure as
  $\SU(2)$, with $C'$ replaced by $C'_N$.
\end{enumerate}
\end{proof}

\begin{corollary}[Scaling with $N$]
The curvature bound~\eqref{eq:kappa-SUN} separates into
$N$-independent and $N$-dependent parts:
\begin{equation}
\kappa_B^{(N)} = \frac{N}{4}
+ \frac{2\beta}{N}\bigl(\cos(\theta_*)\cdot 0.6893
  - C'_2\,\theta_*^2\bigr)
- \beta\,C'_d(N)\,\theta_*^2,
\end{equation}
where $C'_2 \leq 9$ is the $\SU(2)$ non-abelian constant and
$C'_d(N) = O(N)$ is the $d$-symbol channel cost.

The SU(2) corridor condition $\cos\theta_*\cdot 0.6893 > 9\theta_*^2$
(i.e., $\theta_* \lesssim 0.27$) ensures the second term is
positive for \emph{all} $N$. The $d$-symbol cost
$\beta C'_d(N)\theta_*^2$ is the genuinely new obstacle for
$N \geq 3$, but the Ricci term $\frac{N}{4}$ (absent in the
Gaussian approach) compensates: positivity requires
$\frac{N}{4} > \beta C'_d(N)\theta_*^2$, which holds for
any $N$ provided $\beta$ is not too large relative to $N$.
\end{corollary}


%==============================================================================
\section{Consequences for the Curvature Bound}\label{sec:applications}
%==============================================================================

\subsection{Improved corridor for $\SU(2)$}

With the explicit $C'$ from Theorem~\ref{thm:C-prime-su2}:

\begin{corollary}[Corridor width for $\SU(2)$]\label{cor:corridor-su2}
For $\SU(2)$ and $\beta \geq \beta_0$, the curvature bound
$\kappa_B > 0$ holds for all $\theta_*$ satisfying
\begin{equation}
\cos(\theta_*)\cdot 0.6893 > C'\,\theta_*^2.
\end{equation}
With the crude bound $C' \leq 9$: $\theta_* \lesssim 0.27$.
With a tighter bound $C' \leq 5$: $\theta_* \lesssim 0.35$.

The chessboard estimate gives
$\mu(\theta_p > \theta_*) \leq C\exp(-\beta(1-\cos\theta_*)/2)$.
At $\theta_* = 0.27$: $1 - \cos(0.27) \approx 0.0363$, giving
$\mu(\Omega_g^c) \leq Ce^{-0.0182\beta}$. The corridor is
compatible with the chessboard rate for all $\beta \geq \beta_0$.
\end{corollary}


\subsection{Summary: what survives and what cancels}

\begin{center}
\begin{tabular}{llll}
\toprule
Correction & Channel & Order & Quadratic form \\
\midrule
Christoffel ($\Gamma^c_{ab}$)
& $\bigwedge^2(\ad)$
& $O(\theta)$
& $0$ (Theorem~\ref{thm:christoffel}) \\[3pt]
BCH leading ($[A,B]$)
& $\bigwedge^2(\ad)$
& $O(\theta)$
& $0$ (subsumed by Christoffel) \\[3pt]
Scalar metric ($\frac{\sin\theta}{\theta}$)
& $\mathbf{1} \subset \Sym^2$
& $O(\theta^2)$
& $\alpha(\theta)\,\delta_{ab}$ \\[3pt]
Rank-one ($d\theta \otimes d\theta$)
& $V_2 \subset \Sym^2$ ($\SU(2)$)
& $O(\theta^2)$
& Bounded by $C_2\theta_*^2$ \\[3pt]
Jacobian ($\ad^2$)
& $\mathbf{1} \oplus V_2$ ($\SU(2)$)
& $O(\theta^2)$
& Bounded by $C_J\theta_*^2$ \\[3pt]
BCH quadratic ($[A,[A,B]]$)
& $\mathbf{1} \oplus \Sym^2$
& $O(\theta^2)$
& Bounded by $C_B\theta_*^2$ \\[3pt]
$d$-symbol ($N \geq 3$)
& $\ad \subset \Sym^2$
& $O(\theta^2)$
& Bounded by $C_d(N)\theta_*^2$ \\
\bottomrule
\end{tabular}
\end{center}

The key structural result: \textbf{all $O(\theta)$ corrections
are antisymmetric and vanish}. The surviving corrections are all
$O(\theta^2)$, which is what makes the ``wide corridor''
possible.


%==============================================================================
\section{Discussion}\label{sec:discussion}
%==============================================================================

\subsection{What this paper achieves}

We have:
\begin{enumerate}[label=(\roman*)]
\item Promoted the $\SU(2)$ Christoffel vanishing
  of~\cite{Fontes2026Paper8} to a theorem for all compact Lie
  groups with bi-invariant metrics
  (Theorem~\ref{thm:christoffel}).
\item Decomposed the surviving metric correction into
  representation-theoretic channels
  (Theorem~\ref{thm:channel-decomp}), identifying which channels
  contribute to $C'$.
\item Computed an explicit (crude) bound $C' \leq 9$ for
  $\SU(2)$ (Theorem~\ref{thm:C-prime-su2}).
\item Extended the curvature bound to $\SU(N)$
  (Theorem~\ref{thm:curvature-SUN}), identifying the $d$-symbol
  channel as the genuinely new feature for $N \geq 3$.
\end{enumerate}

\subsection{What remains}

\begin{enumerate}[label=(\alph*)]
\item \textbf{Tighter $C'$:} the bound $C' \leq 9$ is crude;
  an exact computation for the $2^4$ block (tracking all 72
  plaquettes and their combinatorial overlaps) would likely give
  $C' \leq 5$. This is a finite computation.
\item \textbf{$d$-channel constant for $\SU(N)$:} the constant
  $C'_d(N)$ bounding the $d$-symbol channel needs explicit
  computation for specific $N$ to determine the effective
  corridor width. The abelian Hessian eigenvalue scales as
  $\lambda_{\min}^{(N)} = \frac{2}{N}\cdot 0.6893$
  (decreasing with $N$), while the $d$-channel correction
  $C'_d(N)$ grows with $N$; the Ricci term $\frac{N}{4}$
  partially compensates.
\item \textbf{The boundary term:} the boundary Poincar\'e
  inequality deferred in~\cite{Fontes2026Paper8} remains open.
\item \textbf{Continuum limit:} uniformity of constants along
  the asymptotically free trajectory $\beta(a) \sim -b_0 \log a$.
\end{enumerate}


\subsection{Relationship to the operator factorization program}

This paper is the ninth in the
series~\cite{Fontes2026a,Fontes2026b,Fontes2026c,Fontes2026i,Fontes2026d,Fontes2026e,Fontes2026f,Fontes2026Paper8}:

\begin{center}
\begin{tabular}{lll}
\toprule
Paper & Domain & Key result \\
\midrule
1~\cite{Fontes2026a} & Hilbert / Gaussian & Factorization identity \\
2~\cite{Fontes2026b} & Rough fBM & $H < 1/2$ extension \\
3~\cite{Fontes2026c} & Non-semimartingale & It\^o formula \\
4~\cite{Fontes2026i} & Energy space / PDE & Generator from $\|\cdot\|_{\cH}$ \\
5~\cite{Fontes2026d} & Banach / stable L\'evy & Leibniz defect \\
6~\cite{Fontes2026e} & Stable--fBM duality & $\Gamma_L = R_H$ \\
7~\cite{Fontes2026f} & Fractional BE & $\kappa(1)=1$ \\
8~\cite{Fontes2026Paper8} & Lattice YM ($\SU(2)$) & Spectral gap \\
\textbf{9 (this paper)} & \textbf{$\SU(N)$ channels}
& \textbf{Christoffel vanishing, $C'$ explicit} \\
\bottomrule
\end{tabular}
\end{center}

The progression from the abstract Leibniz defect
($\Gamma_{\cD}(f,g) = \cD(fg) - f\cD g - g\cD f$) to
representation-theoretic cancellations in $\SU(N)$ gauge theory
demonstrates the unifying power of the operator factorization
framework: the same algebraic structure that produces the
L\'evy--Khintchine generator from jump
defects~\cite{Fontes2026i,Fontes2026d} produces the Wilson
Hessian from non-abelian connection defects, and the cancellation
pattern is governed by the representation theory of the defect's
symmetry group.


\subsection*{Declaration}
AI tools were used as interactive assistants for mathematical
exploration and drafting. The author is responsible for all
content.


%==============================================================================
\begin{thebibliography}{20}

\bibitem{Fontes2026Paper8}
R.~Fontes.
\newblock Non-{G}aussian {B}akry--\'Emery curvature for lattice
  {Y}ang--{M}ills theory via operator factorization.
\newblock \emph{arXiv preprint}, 2026.

\bibitem{Fontes2026a}
R.~Fontes.
\newblock Stochastic calculus as operator factorization.
\newblock \emph{arXiv preprint} arXiv:2601.09976, 2026.

\bibitem{Fontes2026b}
R.~Fontes.
\newblock Stochastic calculus for rough fractional {B}rownian motion
  via operator factorization.
\newblock \emph{arXiv preprint} arXiv:2601.09967, 2026.

\bibitem{Fontes2026c}
R.~Fontes.
\newblock An {I}t\^o formula via predictable projection.
\newblock \emph{arXiv preprint} arXiv:2601.10359, 2026.

\bibitem{Fontes2026i}
R.~Fontes.
\newblock {PDE} generators from energy space structure in
  {F}eynman--{K}ac theory.
\newblock \emph{arXiv preprint}, 2026.

\bibitem{Fontes2026d}
R.~Fontes.
\newblock Operator factorization beyond {H}ilbert spaces: {B}anach
  duality and stable {L}\'evy processes.
\newblock \emph{arXiv preprint}, 2026.

\bibitem{Fontes2026e}
R.~Fontes.
\newblock Stable--fractional {B}rownian duality.
\newblock \emph{arXiv preprint}, 2026.

\bibitem{Fontes2026f}
R.~Fontes.
\newblock Fractional {B}akry--\'Emery curvature on the torus.
\newblock \emph{arXiv preprint}, 2026.

\end{thebibliography}

\end{document}
